% TOP_PROBLEM: {"id": 6600012, "title": "A. Haynes: Gaps Problems — Problem", "category_id": 6, "category": {"id": 6, "name": "geometry", "display_name": "Geometry", "description": "Euclidean and non-Euclidean geometry, geometric structures.", "slug": "geometry", "order_index": 6, "created_at": "2026-07-31T15:26:25.671Z"}, "status": "partially_solved"}
% FIRST_POSTED: null
% ENTRY_KIND: statement_only
% Imported from: batch13.tex; UnsolvedMath ID 6600012
\documentclass[11pt]{article}
\usepackage[margin=1in]{geometry}
\usepackage{amsmath,amssymb,mathtools}
\usepackage{fontspec,unicode-math}
\setmainfont{Latin Modern Roman}
\setsansfont{Latin Modern Sans}
\setmathfont{Latin Modern Math}
\usepackage{xurl,hyperref}
\hypersetup{colorlinks=true,urlcolor=blue,linkcolor=blue}
\newcommand{\sref}[2]{\hyperlink{source:#1:#2}{[S#2]}}
\newcommand{\cataloguescope}[1]{\paragraph{Recorded mathematical question.}#1}
\begin{document}
% UnsolvedMath ID 6600012; source catalogue code AMR-065-0012
\setcounter{section}{6600011}
\section{A. Haynes: Gaps Problems — Problem}\label{6600012}
{\small\sffamily UnsolvedMath ID 6600012 · Geometry}
\subsection{Definitions and mathematical statement}

\paragraph{Shared notation.}$\mathbb N=\{1,2,\ldots\}$, and $\mathbb Z,\mathbb Q,\mathbb R,\mathbb C$ denote the integers, rationals, reals and complex numbers. Any other conventions are specified locally.


Let $E=\{(x,y,\alpha x+\beta y):x,y\in\mathbb R\}\subset\mathbb R^3$, with $1,\alpha,\beta$ linearly independent over $\mathbb Q$. Write $F=\{(0,0,z):z\in\mathbb R\}$ and let $\rho,\rho^*$ be the complementary projections onto $E,F$. The canonical window is $W=\rho^*([0,1]^3)$ and the strip is $S=E+W$. Choose a physical projection $\pi:\mathbb R^3\to E$ whose kernel complements $E$ and whose restriction to $\mathbb Z^3$ is injective. For a regular shift $s$ (no point of $\mathbb Z^3+s$ on $\partial S$), put $Y=\pi(S\cap(\mathbb Z^3+s))$. The canonical system is aperiodic and totally irrational. For $y\in Y$, denote its unique lattice lift by $\widetilde y\in S\cap(\mathbb Z^3+s)$. A type-2 patch of shape $\Omega\subset E$ anchored at $y$ is
\[P_2(y,\Omega)=\{y'\in Y:\rho(\widetilde y'-\widetilde y)\in\Omega\}.\]
Two anchored patches have the same type when $P_2(y,\Omega)-y=P_2(z,\Omega)-z$. Its frequency is the limiting proportion, among anchors $z\in Y\cap B_R$, having that type as $R\to\infty$. Let $\xi_{(\alpha,\beta)}(\Omega)$ be the set of distinct such frequency values, rather than the set of patch types. Aligned rectangles are images in $E$ of coordinate rectangles in the $(x,y)$ plane; aligned squares have equal coordinate side lengths. Thus alignment and squareness use these graph coordinates. These conventions are given in the source’s cited definitions. Put $\|a\|=\min_{m\in\mathbb Z}|a-m|$. The Littlewood condition for $(\alpha,\beta)$ is $\liminf_{n\to\infty}n\|n\alpha\|\|n\beta\|=0$.
\paragraph{Statement recorded in the supplied source.}
For every $(\alpha,\beta)$ with $1,\alpha,\beta$ linearly independent over $\mathbb Q$, does the Littlewood condition imply $\sup_{\Omega\in\mathcal R}\#\xi_{(\alpha,\beta)}(\Omega)=\infty$? \sref{6600012}{2}
\subsection{Short English statement}
Does Littlewood’s approximation condition force unbounded numbers of patch-frequency values over aligned rectangles?
\subsection{Sources}
\begin{description}
\item[\hypertarget{source:6600012:1}{S1}] ULAM.AI and the UnsolvedMath Contributors, \emph{UnsolvedMath: A Curated Collection of Open Mathematics Problems}, record 6600012 (AMR-065-0012). CC BY 4.0. \url{https://huggingface.co/datasets/ulamai/UnsolvedMath}.
\item[\hypertarget{source:6600012:2}{S2}] Faustin Adiceam (compiler), \emph{Open Problems and Conjectures related to the Theory of Mathematical Quasicrystals}, arXiv:1604.06280v2 (2016), Section 2.4, Problem 2.4.3, p. 7. \url{https://arxiv.org/pdf/1604.06280}.
\item[\hypertarget{source:6600012:3}{S3}] Alan Haynes, Henna Koivusalo, Lorenzo Sadun and James Walton, \emph{Gaps problems and frequencies of patches in cut and project sets}, arXiv:1411.0578 (2014), Section 1.2, pp. 2--5. \url{https://arxiv.org/pdf/1411.0578}.
\item[\hypertarget{source:6600012:4}{S4}] Alan Haynes, Henna Koivusalo and James Walton, \emph{Perfectly ordered quasicrystals and the Littlewood conjecture}, arXiv:1506.05649v3 (2015), Sections 1.1 and 2.1, pp. 1--2 and 6--8. \url{https://arxiv.org/pdf/1506.05649}.
\end{description}
\label{end:6600012}
\end{document}
